\documentclass[11pt]{article}
\usepackage[margin=1in]{geometry}
\usepackage[T1]{fontenc}
\usepackage{lmodern,amsmath,amssymb,amsthm,booktabs,array,longtable}
\usepackage[hidelinks]{hyperref}
\newtheorem{theorem}{Theorem}
\newtheorem{lemma}[theorem]{Lemma}
\newtheorem{corollary}[theorem]{Corollary}
\newcommand{\fr}{\operatorname{fr}}
\newcommand{\ioct}{\operatorname{ioct}}
\title{Unbounded healing chains in the question of\\
Choi, Yoo, Kim, Kwon and Won}
\author{Haobo Ma\\\small ChronoAI Pte Ltd, Singapore
\and Wenlin Zhang\\\small National University of Singapore\\
\small ChronoAI Pte Ltd, Singapore}
\date{November 17, 2026}
\begin{document}
\maketitle
\begin{abstract}
The constant-chain-length question of Choi, Yoo, Kim, Kwon and Won
(arXiv:2610.41277v1, Section VIII, Open problems (1), p.~7) has a
negative answer already on outerplanar subcubic graphs of treewidth two:
root-to-leaf forcing makes the required displacement arbitrarily large.
For every full rooted triangle tree, we prove that the minimum displacement
from its nonroot-apex set to an augmentable tight set of the same size
equals its shortest leaf depth. The same number is the minimum number of
tight exchanges. At depth $h$, the smallest graph in this class has
$9\cdot2^h-6$ vertices. Every optimal two-colouring of a tree of triangles
is nevertheless healable, separating the stuck-set obstruction from
failure of a prescribed maximum-cut certificate.
\end{abstract}

\section{The question and the relevant distinctions}

Choi, Yoo, Kim, Kwon and Won ask, ``is the chain length bounded by a
constant?''\ \cite[Section VIII, Open problems (1), p.~7]{CYKKW}
Their Proposition~15 constructs a route out of a stuck tight set with a
bound depending on the number of components of the remaining graph.
We determine the exact escape distance for a class of triangle trees,
including an infinite family in which that distance tends to infinity.

For a finite simple graph $G$, let $\fr(G)$ be the minimum number of
monochromatic edges in a two-colouring of $V(G)$. For $X\subseteq V(G)$,
write $G-X$ for the induced graph on $V(G)\setminus X$ and put
\[
 \varphi_G(X)=|X|+\fr(G-X).
\]
As in Appendix~B of \cite{CYKKW}, a set $X$ is \emph{tight} if it is
independent and $\varphi_G(X)=\fr(G)$. It is \emph{augmentable} if
$X\cup\{v\}$ is tight for some $v\notin X$; a tight set that is not
augmentable is \emph{stuck}. A tight exchange removes one selected vertex
and inserts another, preserving tightness and cardinality. Define
\[
 d_G(S)=\min\{|S\setminus T|: T\text{ augmentable and tight},\ |T|=|S|\}
\]
whenever this set of destinations is nonempty.

The positive identity $\ioct(G)=\fr(G)$ on subcubic graphs without a
$K_4$ component is distinct from this escape question. Here $\ioct(G)$
is the minimum size of an independent odd cycle transversal.
As explained in \cite[Section V and Supplement A.1]{CYKKW}, the identity
follows from the optimal-deletion equalities of Akbari, Beikmohammadi,
Klav\v{z}ar and Movarraei \cite[Theorems 1.1 and 4.1]{ABKM}, whose
preprint appeared in 2021. The constructive chain procedure is the
contribution of \cite[Appendix B]{CYKKW}. Our issue is the cost of
escaping a specified choice, not the existence of an optimal transversal.

There is also a distinction from reconfiguration between two \emph{full}
independent transversals. Buys, Kang and Ozeki prove connectivity under
degree and block-size hypotheses \cite[Theorem 1.2]{BKO}; Hajebi and
Haxell give polynomial-length paths or a dominating-set obstruction
\cite[Theorem 1.4 and Corollary 1.5]{HH}. For triangle blocks below,
deleting the edges within each block leaves a matching, with maximum
degree one and block size three, so their full-transversal setting is
applicable. The nonroot-apex set instead misses the root block, and the
destination ranges over all augmentable sets of that same smaller size.
These full-transversal results do not bound $d_G(S)$ by a constant.

\section{Exact displacement on a rooted triangle tree}

Let $\mathcal T$ be a finite rooted tree whose root has three children
and whose other nodes have either zero or two children. Write $B$ for
its number of nodes and $\delta\geq1$ for its shortest root-to-leaf
distance, measured in edges. Replace each node $u$ by a triangle $Q_u$.
At each nonroot node designate one vertex $a_u$ as its \emph{apex}.
Join each parent triangle to its child's apex by one edge. The root's
three child edges use its three distinct vertices; at any other internal
node, the two child edges use its two nonapex vertices. Call the resulting
graph $G_{\mathcal T}$ and set
\[
 S_{\mathcal T}=\{a_u:u\ne\text{root}\}.
\]
All intertriangle edges are bridges. A tree of this specified form is
called a full rooted triangle tree here.

\begin{theorem}\label{thm:sharp}
The graph $G_{\mathcal T}$ is connected, simple, outerplanar, subcubic,
of treewidth two, and has $3B$ vertices. It is not $K_4$, and
\[
 \fr(G_{\mathcal T})=B,\qquad |S_{\mathcal T}|=B-1,\qquad
 \fr(G_{\mathcal T}-S_{\mathcal T})=1.
\]
The set $S_{\mathcal T}$ is stuck and tight, and
\begin{equation}\label{eq:sharp}
 d_{G_{\mathcal T}}(S_{\mathcal T})=\delta.
\end{equation}
The minimum number of tight exchanges from $S_{\mathcal T}$ to an
augmentable set is also $\delta$. There is an attaining sequence whose
earlier sets are all stuck. Moreover, $B\geq3\cdot2^\delta-2$.
\end{theorem}

\begin{lemma}\label{lem:deletion}
Every independent set $X$ in $G_{\mathcal T}$ satisfies
\[
 \fr(G_{\mathcal T}-X)=B-|X|.
\]
Consequently every independent set is tight.
\end{lemma}
\begin{proof}
An independent set meets each triangle in at most one vertex, so exactly
$B-|X|$ triangles remain intact. Each contributes at least one
monochromatic edge in any two-colouring. For the reverse inequality,
colour each intact triangle with exactly one monochromatic edge and
each remaining two-vertex edge properly. The surviving bridges connect
these pieces according to a forest. Root each component of that forest.
Moving away from its root, flip all colours in a child piece if needed
to make its parent bridge bichromatic. Each piece has at most one
already coloured parent, and flipping preserves its internal cost.
Thus every bridge is bichromatic and the lower bound is attained.
Taking $X=\varnothing$ gives $\fr(G_{\mathcal T})=B$ and hence tightness.
\end{proof}

\begin{proof}[Proof of Theorem~\ref{thm:sharp}]
Each vertex has its two triangle neighbours and at most one bridge
neighbour. The underlying tree gives connectedness and simplicity.
An outerplane drawing is obtained by attaching successive triangles
and their parent bridges in the exterior face, starting at the root;
all vertices remain on that face. For a tree decomposition, take a bag
for each triangle and a two-vertex bag for each bridge, joining a bridge
bag to the two incident triangle bags. The bag graph is a tree, and
the bags containing any fixed vertex are connected. Its maximum bag
size is three, so the width is at most two; the presence of a triangle
forces width at least two. There are $3B\geq12$ vertices, excluding $K_4$.

The set $S=S_{\mathcal T}$ is independent. A bridge has a nonapex
parent endpoint, except at the root, where no vertex belongs to $S$.
It has size $B-1$, and deleting it leaves the root triangle and
$B-1$ isolated edges. Lemma~\ref{lem:deletion} gives tightness and
residual frustration one. Every vertex outside $S$ has a neighbour in
$S$: a root vertex sees its child's apex, and a nonroot nonapex vertex
sees its own apex. Thus $S$ is stuck.

For the lower bound, let $T$ be augmentable and tight with $|T|=B-1$,
and choose $v$ such that $R=T\cup\{v\}$ is independent and tight.
Since $|R|=B$, it selects exactly one vertex per triangle. The selected
root vertex is incident to the bridge to some child $u_1$, so
$a_{u_1}\notin R$. The selected vertex of $Q_{u_1}$ is therefore
nonapex. If $u_1$ is internal, that vertex points to a child $u_2$,
again excluding $a_{u_2}$. Continue to a leaf at depth $\ell\geq\delta$.
The distinct apices on this path lie in
$S\setminus R\subseteq S\setminus T$. Hence $|S\setminus T|\geq\delta$.

To attain the bound, choose a shortest root-to-leaf path
$u_0,u_1,\ldots,u_\delta$. For $0\leq i<\delta$, let $z_i$ be the
parent endpoint of the bridge from $Q_{u_i}$ to $a_{u_{i+1}}$; let
$z_\delta$ be either nonapex vertex of the leaf triangle. Set
\[
 R=\bigl(S\setminus\{a_{u_1},\ldots,a_{u_\delta}\}\bigr)
       \cup\{z_0,z_1,\ldots,z_\delta\}.
\]
This selects one vertex in each triangle. On a bridge between path
triangles the child apex is absent; on a bridge leaving the path the
parent endpoint is absent. Bridges wholly off the path also have their
nonapex parent endpoints absent. Thus $R$ is independent.
Put $T=R\setminus\{z_\delta\}$. Both sets are tight by
Lemma~\ref{lem:deletion}; $T$ is augmentable and misses exactly the
$\delta$ displayed apices of $S$, proving \eqref{eq:sharp}.

For $0\leq j\leq\delta$ put
\[
 T_j=\bigl(S\setminus\{a_{u_1},\ldots,a_{u_j}\}\bigr)
          \cup\{z_0,\ldots,z_{j-1}\},
\]
with both displayed sets empty at $j=0$. These are independent subsets
of $R$, hence tight, and each step exchanges $a_{u_{j+1}}$ for $z_j$.
The displacement of $T_j$ is $j$, so the lower bound makes every
$T_j$ with $j<\delta$ stuck, whereas $T_\delta=T$ is augmentable.
Any sequence of $m$ single-vertex exchanges has displacement at most $m$;
therefore this sequence is shortest. Finally all nodes at depths less
than $\delta$ are internal, so
$B\geq1+3\sum_{i=0}^{\delta-1}2^i=3\cdot2^\delta-2$.
\end{proof}

\begin{corollary}\label{cor:uniform}
For every $h\geq1$, take all leaves at depth $h$. The resulting graph
$G_h$ has $B=3\cdot2^h-2$ triangles and $n=9\cdot2^h-6$ vertices,
and its stuck apex set has exact displacement $h$. Thus no constant
bound holds even for outerplanar graphs of treewidth two.
Among full rooted triangle trees with shortest leaf depth $h$, this
order is the smallest possible. More generally,
\[
 d_{G_{\mathcal T}}(S_{\mathcal T})
 \leq\left\lfloor\log_2\frac{|V(G_{\mathcal T})|+6}{9}\right\rfloor,
\]
with equality on $G_h$.
\end{corollary}
\begin{proof}
Count the three branches of the root and the binary levels below them.
Theorem~\ref{thm:sharp} gives the exact distance and the matching
order lower bound. Solving that lower bound for $\delta$ gives the
displayed inequality. This is an upper bound within this triangle-tree
class, not an upper bound for all subcubic graphs.
\end{proof}

\section{Chain stages and the prescribed-colouring sibling}

\begin{corollary}\label{cor:stages}
For $G_{\mathcal T}$ and $S_{\mathcal T}$, the smallest possible
zero-based stage index of a successful test in Proposition~15 of
\cite{CYKKW} is $\delta-1$. In particular, $G_h$ requires at least
$h$ tested stages, counting the successful stage.
\end{corollary}
\begin{proof}
At stage $j$, the maintained set has undergone $j$ exchanges from $S$.
One further test exchange gives a set $U$ with
$|S\setminus U|\leq j+1$. If $U$ is augmentable,
Theorem~\ref{thm:sharp} implies $j+1\geq\delta$.
The sequence $T_j$ in its proof is a chain of the source's kind:
$G_{\mathcal T}-T_j$ has $Q_{u_j}$ as its sole intact triangle,
with all other components isolated edges. Exchanging
$a_{u_{j+1}}$ for $z_j$ moves the odd triangle to $Q_{u_{j+1}}$.
The first successful test is therefore at $j=\delta-1$.
\end{proof}

There are $B$ components in $G_{\mathcal T}-S_{\mathcal T}$, so the
source's component-dependent upper bound is retained. Its observation
that the tests used at most three stages is a finite experimental
statement. A universal three-stage bound is not asserted in the source;
it fails for $G_4$ (138 vertices). The polynomial construction of
\cite[Corollary 16]{CYKKW} is consistent with these logarithmic
obstructions.

The source also distinguishes a healable Ising ground state from a
graph on which every Ising ground state is healable
\cite[Section V and Appendix A, ``Strong version'']{CYKKW}.
A two-colouring is \emph{healable} if one can choose exactly one
endpoint of each monochromatic edge so that the chosen vertices are
independent. The following restricted positive result shows why the
escape obstruction is different.

\begin{corollary}\label{cor:healable}
Let $G$ be obtained from vertex-disjoint triangles by adding bridges
whose contraction gives a tree. Every optimal two-colouring of $G$ is
healable. No degree restriction is needed for this assertion.
\end{corollary}
\begin{proof}
The colouring argument of Lemma~\ref{lem:deletion} with $X=\varnothing$
works for this more general bridge pattern, so $\fr(G)=B$, the number
of triangles. Every optimal colouring therefore has exactly one
monochromatic edge per triangle and none on the bridges.
Root the contracted tree. In the root triangle choose either endpoint
of its monochromatic edge. Proceed away from the root. The selected
vertex of a triangle's parent can forbid at most one of its two
candidate endpoints, because exactly one bridge joins them; choose a
remaining endpoint. Every bridge constraint is met when its child is
processed, so the resulting $B$ choices are independent. They contain
exactly one endpoint of every monochromatic edge.
\end{proof}

\section{Coverage and verification}

\renewcommand{\arraystretch}{1.12}
\begin{longtable}{@{}p{0.32\linewidth}p{0.63\linewidth}@{}}
\toprule
Source item & Mathematical coverage \\
\midrule
\endfirsthead
\toprule
Source item & Mathematical coverage \\
\midrule
\endhead
\cite{CYKKW}, Section VIII, question (1), p.~7 &
Answered: no. Corollaries~\ref{cor:uniform} and \ref{cor:stages}
give unbounded exact distances and stage counts. \\
Uneven full triangle trees &
Not posed; proved. Theorem~\ref{thm:sharp} identifies shortest leaf
depth as the exact distance; Corollary~\ref{cor:uniform} gives the
sharp order bound within this class. \\
\cite{CYKKW}, Section VIII, three-stage observation &
Finite observation retained. A universal reading is not asserted in
the source; it fails for $G_4$. \\
\cite{CYKKW}, Proposition 15, pp.~10--11; Corollary 16, p.~11 &
Prior constructive results cited. Their component-dependent and
polynomial bounds are unchanged. \\
\cite{CYKKW}, Appendix A, stuck-set examples, p.~9; Appendix B, p.~11 &
Prior 12-vertex example cited. Theorem~\ref{thm:sharp} proves an
infinite family and its exact escape distances. \\
\cite{CYKKW}, Appendix A, strong version, pp.~8--9 &
Prior general counterexample cited. The restricted assertion in
Corollary~\ref{cor:healable} is not posed there; proved for every
bridge-connected tree of triangles. \\
\cite{CYKKW}, Theorem 8 and Supplement A.1 &
Prior positive transversal identity cited, with its antecedent
\cite[Theorems 1.1 and 4.1]{ABKM}. \\
\cite{ABKM}, Problem 3.2, pp.~7--8 &
Prior answer cited: no, by Cambie, Haslegrave and Kang
\cite{CHK}; planar degree-four examples are in
\cite[Theorem 1.2]{KPS}. Not a new claim here. \\
\cite{CYKKW}, Section VIII, questions (2)--(4), pp.~7--8 &
Not covered: the higher-degree complexity window, Penrose
thermodynamic kinks, and finite-temperature excitations. \\
Other energy, hardness and reduction statements of \cite{CYKKW},
including Supplement A.2--A.7 &
Not covered by the triangle-tree argument. \\
\bottomrule
\end{longtable}

The pinned
\href{https://github.com/the-omega-institute/trureturing/blob/0ad94071772b2f32b432905d088bfce029d8d29e/D5/S3/Combinatorics/Graph/Potts/HealingUnbounded.lean}{\texttt{HealingUnbounded.lean}}
contains the declaration
\begin{quote}\small
\nolinkurl{D5.S3.Combinatorics.Graph.Potts.HealingUnbounded.result}.
\end{quote}
Its statement is the negation of \texttt{claim}: no natural number $C$
uniformly bounds displacement to an augmentable tight set of equal size
over finite connected subcubic graphs other than $K_4$, starting from
stuck tight sets with positive residual frustration. Its definitions
of frustration, induced vertex deletion, tightness and augmentation
agree with those above. Its construction is $G_h$, indexed by root
block $0$, root children $1,2,3$, and children $2u+2,2u+3$ at each
positive internal block $u$.

The exact shortest-leaf formula, attaining sequences, outerplanarity,
treewidth, deletion identity, stage-index interpretation and universal
healability within triangle trees have written proofs in this note.
The accompanying exact-integer programs give separate finite checks:
\texttt{verify.py} covers uniform heights $1$--$10$ (up to 9,210
vertices), including small exhaustive enumerations;
\texttt{verify\_siblings.py} covers all 428 ordered full rooted trees
with $4$--$14$ triangles and all 162 optimal two-colourings of $G_1$.
The ranges and outputs are in \texttt{README.md}, \texttt{checks.json}
and \texttt{sibling-checks.json}. These finite checks supplement the
proofs for arbitrary finite trees.

\begin{thebibliography}{9}
\bibitem{CYKKW}
J. Choi, S. Yoo, C. Kim, H. Y. Kwon and C. Won,
\emph{Coordination three heals, coordination four is hard: the
antiferromagnetic Potts model with a costly third state},
\href{https://arxiv.org/abs/2610.41277v1}{arXiv:2610.41277v1}
(October 30, 2026), Section VIII, pp.~7--8; Appendix B, pp.~9--11;
Supplemental Material, A.1.

\bibitem{ABKM}
S. Akbari, A. Beikmohammadi, S. Klav\v{z}ar and N. Movarraei,
\emph{On the chromatic vertex stability number of graphs},
European Journal of Combinatorics \textbf{102}, 103504 (2022).
\href{https://doi.org/10.1016/j.ejc.2021.103504}{doi:10.1016/j.ejc.2021.103504}.
\href{https://arxiv.org/abs/2108.12994v2}{arXiv:2108.12994v2},
Theorems 1.1 (p.~2) and 4.1 (p.~8).

\bibitem{BKO}
P. Buys, R. J. Kang and K. Ozeki,
\emph{Reconfiguration of independent transversals},
Random Structures \& Algorithms \textbf{67}, e70025 (2025).
\href{https://doi.org/10.1002/rsa.70025}{doi:10.1002/rsa.70025}.
\href{https://arxiv.org/abs/2407.04367v1}{arXiv:2407.04367v1}
(July 5, 2024), Theorem 1.2.

\bibitem{HH}
S. Hajebi and P. Haxell,
\emph{Efficient reconfiguration of independent transversals},
\href{https://arxiv.org/abs/2611.18474v1}{arXiv:2611.18474v1}
(November 13, 2026), Theorem 1.4 (p.~3) and Corollary 1.5 (p.~4).

\bibitem{CHK}
S. Cambie, J. Haslegrave and R. J. Kang,
\emph{When removing an independent set is optimal for reducing the
chromatic number}, European Journal of Combinatorics \textbf{115},
103781 (2024).
\href{https://doi.org/10.1016/j.ejc.2023.103781}{doi:10.1016/j.ejc.2023.103781}.
\href{https://arxiv.org/abs/2203.13833v2}{arXiv:2203.13833v2},
Theorem 2; first preprint March 25, 2022.

\bibitem{KPS}
M. Knor, M. Petru\v{s}evski and R. \v{S}krekovski,
\emph{On chromatic vertex stability of 3-chromatic graphs with
maximum degree 4}, Discrete Mathematics Letters \textbf{10},
75--80 (2022), Theorem 1.2 (p.~76).
\href{https://doi.org/10.47443/dml.2022.066}{doi:10.47443/dml.2022.066}.
\end{thebibliography}

\medskip
\noindent\small
\textcopyright\ 2026 Haobo Ma and Wenlin Zhang.
Licensed under Creative Commons Attribution 4.0 International (CC BY 4.0).
\end{document}
